\subsection{有限长度模的结构}

定理 2（克鲁尔–雷马克–施密特）. —— 设 A 是一个环，M 是一个有限长度的 A-模。

\begin{enumerate}
\def\labelenumi{\alph{enumi})}
\item
  存在 M 的不可分解子模的一个有限族 \((\mathrm{M}_i)_{i\in \mathrm{I}}\)，使 \(\mathrm{M} = \bigoplus_{i\in \mathrm{I}}\mathrm{M}_i\)，且模 M 是半原始的。
\item
  设 \((\mathrm{M}_{i})_{i\in\mathrm{I}}\) 与 \((\mathrm{M}_{j}^{\prime})_{j\in\mathrm{J}}\) 是 M 的不可分解子模的两个有限族，使得 \(M=\bigoplus_{i\in I}M_{i}=\bigoplus_{j\in J}M_{j}^{\prime}\) 。存在从 I 到 J 的一个双射 \(\sigma\) 与 M 的一个自同构 u，使 \(u(\mathrm{M}_{i})=\mathrm{M}_{\sigma(i)}^{\prime}\) 对每个 \(i\in I\) 成立。
\item
  设 \(\mathrm{N}\) 是 \(\mathrm{M}\) 的一个直因子子模，并设 \((\mathrm{M}_i)_{i\in \mathrm{I}}\) 是 \(\mathrm{M}\) 的不可分解子模的一个有限族，其直和为 \(\mathrm{M}\) 。存在子集 \(\mathrm{J}\) ⊂ \(\mathrm{I}\)，使 \(\bigoplus_{i\in \mathrm{I - J}}\mathrm{M}_i\) 与 \(\mathrm{N}\) 互补。模 \(\mathrm{N}\) 同构于 \(\bigoplus_{j\in \mathrm{J}}\mathrm{M}_j\) 。
\item
  设 \(\mathbf{N}\) 是一个 A-模。若存在整数 \(d > 0\) 使模 \(\mathbf{M}^d\) 与 \(\mathbf{N}^d\) 同构，则模 \(\mathbf{M}\) 与 \(\mathbf{N}\) 同构。
\item
  设 \(\mathrm{N}\) 与 \(\mathrm{P}\) 是有限长度的 A-模。若模 \(\mathrm{M} \oplus \mathrm{P}\) 与 \(\mathrm{N} \oplus \mathrm{P}\) 同构，则 \(\mathrm{M}\) 与 \(\mathrm{N}\) 同构。
\end{enumerate}

有限长度的模既是阿廷的又是诺特的（VIII，第 2 页，命题 1）。此外，对有限长度的模而言，不可分解与原始是一回事（VIII，第 31 页，命题 4）。于是断言 a) 由 VIII，第 30 页的命题 3 得出。断言 b)、c)、e) 分别由 VIII，第 32 页定理 1 的推论 2、推论 5、推论 3 得出。最后，断言 d) 由 VIII，第 35 页的推论 4 得出，因为在 d) 的假设下，模 N 长度有限，从而由 a) 是半原始的。

定理 3. —— 设 K 是一个交换域，A 是一个 K-代数，M 与 N 是有限长度的模。设 \(K'\) 是一个非零交换 K-代数，使得 \(\mathrm{A}_{(\mathrm{K}')}\) -模 \(\mathrm{M}_{(\mathrm{K}')}\) 与 \(\mathrm{N}_{(\mathrm{K}')}\) 同构。则 A-模 M 与 N 同构。

\begin{enumerate}
\def\labelenumi{\alph{enumi})}
\item
  先设代数 \(K'\) 在 K 上具有有限次数 d。则 A-模 \(\mathrm{M}_{(\mathrm{K}')}\) 同构于 \(M^{d}\)，A-模 \(\mathrm{N}_{(\mathrm{K}')}\) 同构于 \(\mathbf{N}^d\)，故 A-模 \(\mathbf{M}^d\) 与 \(\mathbf{N}^d\) 同构。由定理 2, d)，A-模 \(\mathbf{M}\) 与 \(\mathbf{N}\) 同构。
\item
  现在设 K-代数 \(K'\) 由有限多个元素生成。取 \(K'\) 的一个极大理想 m，并令 \(K'' = K'/m\) 。由希尔伯特零点定理（VIII，第 462 页，定理 1 的推论 1），\(K''\) 是 K 的一个有限次扩张。把标量从 \(K'\) 扩张到 \(K''\) ，我们由 \(A_{(K')}\) -线性同构 \(M_{(K')} \to N_{(K')}\) 得到 \(A_{(K'')}\) -线性同构 \(M_{(K'')} \to N_{(K'')}\) 。由证明的 a)，A-模 M 与 N 同构。
\item
  最后处理一般情形。设 \(u: \mathrm{M}_{(\mathrm{K}')} \to \mathrm{N}_{(\mathrm{K}')}\) 是 \(\mathrm{A}_{(\mathrm{K}')}\) -模的一个同构，\(v: \mathrm{N}_{(\mathrm{K}')} \to \mathrm{M}_{(\mathrm{K}')}\) 是其逆同构。用 \(\mathcal{E}\) 表示由有限多个元素生成的 \(\mathrm{K}'\) 的 K-子代数组成的集合。若 E 是这样的一个子代数，则 \(\mathrm{A}_{(\mathrm{E})}\) 等同于 \(\mathrm{A}_{(\mathrm{K}')}\) 的一个子环，且 \(\mathrm{M}_{(\mathrm{E})}\) 与 \(\mathrm{N}_{(\mathrm{E})}\) 分别等同于 \(\mathrm{A}_{(\mathrm{E})}\) -子模——前者在 \(\mathrm{M}_{(\mathrm{K}')}\) 中，后者在 \(\mathrm{N}_{(\mathrm{K}')}\) 中（II, §7, No.~7, 第 306 页）；此外，\(\mathrm{M}_{(\mathrm{K}')}\) 与 \(\mathrm{N}_{(\mathrm{K}')}\) 分别是右有向族 \((\mathrm{M}_{(\mathrm{E})})_{\mathrm{E} \in \mathcal{E}}\) 与 \((\mathrm{N}_{(\mathrm{E})})_{\mathrm{E} \in \mathcal{E}}\) 的并。A-模 M 与 N 长度有限，从而是有限生成的；设 S 是 A-模 M 的一个有限生成元集，T 是 A-模 N 的一个有限生成元集。存在 K-代数 \(\mathrm{E} \in \mathcal{E}\)，使 \(u(1 \otimes s) \in \mathrm{N}_{(\mathrm{E})}\) 对每个 \(s \in \mathrm{S}\) 成立，且 \(v(1 \otimes t) \in \mathrm{M}_{(\mathrm{E})}\) 对每个 \(t \in \mathrm{T}\) 成立。由线性性推出 \(u(\mathrm{M}_{(\mathrm{E})}) \subset \mathrm{N}_{(\mathrm{E})}\) 与 \(v(\mathrm{N}_{(\mathrm{E})}) \subset \mathrm{M}_{(\mathrm{E})}\) 。于是映射 \(u\) 与 \(v\) 诱导出从 \(\mathrm{M}_{(\mathrm{E})}\) 到 \(\mathrm{N}_{(\mathrm{E})}\) 以及从 \(\mathrm{N}_{(\mathrm{E})}\) 到 \(\mathrm{M}_{(\mathrm{E})}\) 的互逆双射。这些双射显然是 \(\mathrm{A}_{(\mathrm{E})}\) -线性的。于是 \(\mathrm{A}_{(\mathrm{E})}\) -模 \(\mathrm{M}_{(\mathrm{E})}\) 与 \(\mathrm{N}_{(\mathrm{E})}\) 同构。由证明的 b)，A-模 M 与 N 同构。
\end{enumerate}

注. —— 设 E 与 F 是交换域 K 上的两个有限维向量空间，并设 \(K'\) 是 K 的一个扩张。设 u 是 E 的一个自同态，v 是 F 的一个自同态，并设经标量扩张得到的 \(u_{(K')}\) 与 \(v_{(K')}\) 是 \(E_{(K')}\) 与 \(F_{(K')}\) 的自同态。由 VII, §5, No.~3, 第 32 页的推论 1 与推论 2，自同态 u 与 v 相似当且仅当自同态 \(u_{(K')}\) 与 \(v_{(K')}\) 相似。这一点也可由把上面的定理 3 应用于代数 A = K{[}X{]} 及 A-模 \(M = E_u\) 与 \(N = F_v\)（VII, §5, No.~1, 第 29 页）而得到。

